\Th Пусть $K \supset F$ -- конечное расширение, $char \ F = 0$. Тогда найдется примитивный элемент $\gamma \in K$, т.е. такой, что $K = F(\gamma)$.

\Proof
Так как $K = F(\alpha_1, \ldots, \alpha_n)$, то достаточно доказать утверждение о том, что $F(\alpha, \beta)$ порождается над $F$ одним элементом $\gamma$. Обозначим через $f(x)$ (соотв. $g(x)$) минимальный многочлен над $F$ элемента $\alpha$ (соотв. $\beta$).\\

Идея доказательства в следующем: мы докажем, что $\gamma$ можно взять в виде $\gamma = \alpha + c\beta$ при подходящем
$c \in F$. Окажется, что «плохие» $c$ образуют лишь конечное множество, а значит, ввиду бесконечности $F$ такое $c$ всегда найдётся.\\

Давайте найдем условие на $c$, при котором $F(\alpha, \beta) = F(\gamma)$. Для этого достаточно показать, что $\beta \in F(\gamma)$. Действительно, тогда $\alpha = \gamma - c\beta$ также лежит в $F(\gamma)$, а, значит, $F(\alpha, \beta) \subset F(\gamma)$, и из обратного включения мы получаем требуемое равенство.

Рассмотрим многочлен $h(x) = f(\gamma - cx) \in F(\gamma)[x]$. Элемент $\beta$ является корнем этого многочлена.

Если $deg(h, g) = 1$, то $\beta$ является корнем многочлена первой степени, то есть лежит в $F(\gamma)$.\\

Посмотрим, когда $deg(h, g) > 1$. Так как $g$ не имеет кратных корней, то в $\overline{F}$ найдется корень $\beta_1$ многочлена $g(x)$, отличный от $\beta$, который  является корнем многочлена $h$. Тогда $\alpha_1 = \gamma - c\beta_1$ -- корень многочлена $f(x)$. 

Тогда $\alpha + c\beta = \gamma = \alpha_1 + c\beta_1$, и $c = \frac{\alpha_1 - \alpha}{\beta_1 - \beta}$.

Следовательно, если взять $c \neq \frac{\alpha_1 - \alpha}{\beta_1 - \beta}$ ни для каких сопряженных с $\alpha$ и $\beta$ элементов $\alpha_1$ и $\beta_1$, то $\beta$ будет единственным общим корнем $h, g$, а, значит, $\beta \in F(\gamma)$
\EndProof